% ===============================================================
\section{Conclusions}
\label{sec:conclusions}
% ===============================================================

\subsection{Effect of the initial guesses on the stability and speed of convergence}
About 300 GA runs were made in this project (Tables~\ref{tab:t22study}, \ref{tab:t23study},
\ref{tab:t25study} and~\ref{tab:t25ngen}). Table~\ref{tab:concl} summarises how the initial guesses
of the constants affect the algorithm.

\begin{table}[H]
\centering\footnotesize
\caption{Summary: how the initial guesses affect the stability and the speed of convergence of the genetic
algorithm (MFRAC $=0.5$ and MUT $\approx\pm9\,\%$ unless stated).}
\label{tab:concl}
\begin{tabularx}{\linewidth}{@{}>{\raggedright\arraybackslash}p{3.3cm} >{\raggedright\arraybackslash}X >{\raggedright\arraybackslash}X >{\raggedright\arraybackslash}p{3.6cm}@{}}
\toprule
Initial guesses & Observed behaviour & Cause & Remedy\\
\midrule
Close to the solution (predicted output of the right order of magnitude): 2.2 Sets 1--2, 2.3 Set C, 2.5 Sets 3--4
& Stable. Target reached in \textbf{2--12 generations} in every run; final error within a few percent of the true minimum.
& Only the level of the prediction ($n_1$) has to be adjusted.
& Continue with small mutations to refine.\\
\addlinespace
Right order of magnitude but wrong position along a valley: 2.2 Set 1, 2.3 Set A
& Target reached quickly, but the constants depend on the guess and on the seed (e.g.\ $n_2=3.08$--$3.34$) and drift only slowly.
& Correlated constants ($\ln n_1$ and $n_2$; $n_3$, $n_4$ and $n_5$): long, flat valley of the error.
& Several seeds and guesses; larger MUT, then refinement; check the constants, not only the error.\\
\addlinespace
Far from the solution: 2.2 Sets 3--4, 2.3 Set D
& \textbf{Stagnation} at 2--4 times the minimum error. NGEN $=20000$ does not help. With MUT $=\pm40\,\%$: convergence after $\sim$130--1600 generations.
& Steps of at most $\pm$MUT per generation; most mutations away from the valley floor are rejected.
& Two stages: MUT $=0.2$--$0.4$ to explore, then MUT $=0.03$--$0.01$ to refine.\\
\addlinespace
All organisms with nearly the same large error: 2.3 Set E, 2.5 Sets 1--2 $(1,1,1)$, $(2,2,0.5)$
& \textbf{Population frozen}: nkeep $=0$ at every generation, error constant (0.6--0.95). Unaffected by MUT and NGEN.
& Relative threshold $0.5\times$median selects nobody; all offspring come from \code{n[0]}, which is never replaced.
& MFRAC $=1$ (converges in 55--84 generations), or better guesses.\\
\addlinespace
A constant guessed close to zero: 2.3 Set B ($n_5=0.01$)
& The constant stays near zero; only 3/10 runs reach the target.
& Multiplicative mutation: the step is proportional to the value.
& Guess a non-zero value of the expected order of magnitude.\\
\addlinespace
Constants not identifiable from the data: 2.5 ($(UA)_e$, $(UA)_c$, $\alpha$)
& Same error for $(UA)_e$ from 5 to 11~W/\textdegree C; the final constants follow the guess.
& The data only constrain the series conductance $UA_{eff}$.
& Report $UA_{eff}$; physically motivated guess $(UA)_e\approx(UA)_c$; more measurements.\\
\bottomrule
\end{tabularx}
\end{table}

In short:
\begin{enumerate}
  \item \textbf{The most important property of the initial guesses is that the model output they
        predict has the right order of magnitude.} When it does, the GA is stable and fast, whatever the
        exact values. When all organisms are far from the data, the relative selection threshold
        (MFRAC$\,\times\,$median) stops working and the skeleton algorithm can freeze completely.
  \item \textbf{Speed of convergence}: with good guesses, the target error is reached in a few to a few
        tens of generations, and the minimum is found within a few thousand generations. Beyond that,
        increasing NGEN does not help a trapped population. What helps is a larger mutation amplitude
        (exploration) followed by a small one (refinement), or a larger MFRAC.
  \item \textbf{Reaching the target error is not proof of an accurate fit.} With correlated constants,
        many sets of constants meet the target. To be confident, several seeds and initial guesses must give
        the same constants. The deterministic check showed that my best GA fits are within
        $0.03\,\%$ (Task~2.2), $0.6\,\%$ (Task~2.3) and $0.01\,\%$ (Task~2.5) of the true minimum error.
  \item \textbf{Constants that the data do not determine} (Task~2.5) stay at values that depend on the
        initial guess. Only identifiable combinations should be interpreted.
\end{enumerate}

\subsection{Main results}
\begin{itemize}
  \item \textbf{Task 2.2}: $\qpp=\TwoTwoNone\,\dTsat^{\TwoTwoNtwo}g^{\TwoTwoNthree}$, $\Ferrmean=\TwoTwoF$.
        The rms deviation ($\TwoTwoRmsRel\,\%$) is larger than the $\pm10\,\%$ uncertainty of the data, and
        the deviations are systematic between gravity levels. A power law in $g$ cannot represent the
        Marangoni effect.
  \item \textbf{Task 2.3}: $\qpp=\TwoThreeNone\,\dTsat^{\TwoThreeNtwo}(g+\TwoThreeNfour\,\gen\gamma)^{\TwoThreeNthree}P^{\TwoThreeNfive}$,
        $\Ferrmean=\TwoThreeF$ (validation $\TwoThreeFval$). The overall deviation is close to the noise
        level, but the $1\,\gen$ and $2\,\gen$ (5.5~kPa) groups keep a bias of $-16\,\%$ and $+11\,\%$. The
        Marangoni effect acts like an additional gravity of about $2.4\,\gen$ for $\gamma=1.79$, which
        explains why the heat flux hardly depends on gravity for this mixture.
  \item \textbf{ET1}: the heat flux is about 6 times more sensitive to the superheat than to pressure, the
        next most influential input. Keeping $|\Delta\qpp/\qpp|<0.15$ requires about $3.6\,\%$ accuracy on all inputs, or
        about $\pm1.4$--$1.9$\,\textdegree C on the superheat.
  \item \textbf{ET2}: the directly fitted inverse relation differs from the inverted forward relation by
        $\approx5$--$10\,\%$. This difference comes from the scatter of the data (regression dilution),
        as shown with synthetic data.
  \item \textbf{Tasks 2.4--2.5}: $(UA)_e=\TwoFiveNone$~W/\textdegree C, $(UA)_c=\TwoFiveNtwo$~W/\textdegree C, $\alpha=\TwoFiveNthree$.
        The MAE is $\TwoFiveFtrain$ for the training set and $\TwoFiveFval$ for the validation set, well
        below 0.045. Only the overall conductance $UA_{eff}=\TwoFiveUA$~W/\textdegree C is determined uniquely by the data.
\end{itemize}
